(1) 能量守恒与傅里叶定律给 \(\partial_tT=\kappa\partial_x^2T\)，\(\kappa=k/(\rho c)\)。条件为 \(T(x,0)=T_0\)（\(x>0\)）、\(T(0,t)=T_s\)（\(t>0\)）、\(T(\infty,t)=T_0\)。初始边界在 \((x,t)=(0,0)\) 形成温度不连续，因此不能强求此点满足光滑解。
(2) 取 \(\eta=x/[2\sqrt{\kappa t}]\) 和 \(T=T_0+(T_s-T_0)f(\eta)\)，代入扩散方程得到 \(f''+2\eta f'=0\)。利用 \(f(0)=1\)、\(f(\infty)=0\)，解得 \(f(\eta)=\operatorname{erfc}\eta\)，故 \(T(x,t)=T_0+(T_s-T_0)\operatorname{erfc}[x/(2\sqrt{\kappa t})]\)。对任意固定 \(x>0\)，\(t\to0^+\) 时 \(\eta\to\infty\)，确实回到初温。
(3) 由 \(\operatorname{erfc}'(\eta)=-2e^{-\eta^2}/\sqrt\pi\)，表面热流 \(j(0,t)=-k\partial_xT|_{0^+}=k(T_s-T_0)/\sqrt{\pi\kappa t}\)。积分得到 \(Q(t)/A=\int_0^tj(0,s)\,ds=2k(T_s-T_0)\sqrt{t/(\pi\kappa)}\)。这两个表达式分别按 \(t^{-1/2}\) 衰减与 \(t^{1/2}\) 增长，说明深处热扩散尺度随 \(\sqrt{\kappa t}\) 展开。
(4) 内部储热增量 \(Q_{\rm stored}/A=\rho c\int_0^\infty[T(x,t)-T_0]\,dx\)。用 \(\int_0^\infty\operatorname{erfc}\eta\,d\eta=1/\sqrt\pi\)，得到 \(Q_{\rm stored}/A=2\rho c(T_s-T_0)\sqrt{\kappa t/\pi}=Q(t)/A\)，与第三问守恒。\(t\to0^+\) 时理想阶跃边界导致瞬时热流无界，但积分热量仍有限；真实界面接触热阻及热弛豫会钝化此奇性。
